\section{Анализ предметной области}

\subsection{Характеристика предметной области}
% Описываются границы предметной области, ключевые объекты, процессы, ограничения и участники. Следует зафиксировать терминологию, используемую далее в пояснительной записке.
Границами предметной области проекта являются процессы накопления, систематизации, семантического анализа и извлечения знаний из слабоструктурированных и неструктурированных текстовых документов: научные статьи, техническая документация, конспекты, аналитические отчеты. 

В рамках настоящей работы под данными понимаются необработанные текстовые файлы; под информацией — структурированные метаданные (карточки знаний), извлечённые из данных; под знанием — связная система концепций и отношений, обладающая свойством интерпретируемости и пригодная для вывода. Такое разграничение соответствует классической пирамиде DIKW (Data–Information–Knowledge–Wisdom) и определяет место проектируемой системы на уровне «знание».

Ключевыми объектами предметной области выступают:
\begin{itemize}
    \item \textbf{оригинал документа} --- текстовый или бинарный файл (PDF, Markdown, docx), сохраняемый в неизменном виде во внешнем хранилище данных и являющийся первичным источником знаний;
    \item \textbf{карточка знания} --- структурированный объект метаданных, содержащий краткое аннотированное резюме (summary), ключевые тезисы, список связанных концепций и прямую ссылку на оригинал документа;
    \item \textbf{концепция} --- базовая семантическая единица знания, описывающая конкретный термин, метод или объект предметной области, представленная в графе в виде узла;
    \item \textbf{канонический реестр концепций} --- централизованный словарь терминов и их синонимов, обеспечивающий единую идентификацию концепций и предотвращающий дублирование элементов знаний;
    \item \textbf{домен и сообщество} --- укрупненные иерархические уровни группировки концепций, формируемые алгоритмами топологической кластеризации для обеспечения контекстной навигации;
    \item \textbf{семантический граф знаний} --- связный ориентированный граф $G = (V, E)$, где множеством вершин $V$ являются концепции и домены, а множеством рёбер $E$ --- смысловые и ассоциативные связи между ними, наделенные весовыми коэффициентами.
\end{itemize}

Система оперирует тремя видами знаний: декларативными (факты о концепциях и документах), структурными (иерархия «сообщество → домен → концепция») и реляционными (семантические связи между концепциями с весами). Процедурные знания (правила вывода) в текущей версии не формализуются, однако архитектура допускает их добавление на уровне агентов.

Участниками процесса являются пользователи системы:аналитики, исследователи, инженеры и т.д., загружающие новые документы и формулирующие поисковые или аналитические запросы на естественном языке, а также автономные программные агенты, выполняющие специализированные функции обработки данных.


К основным ограничениям предметной области относятся: динамический характер поступающих знаний (структура базы не является фиксированной), недопустимость избыточных вычислительных затрат уровня $O(V+E)$ при обработке запроса в рантайме, а также необходимость работы системы в условиях высокой синонимии и терминологической неоднородности исходных текстов.

Дополнительным ограничением является омонимия — одна и та же лексема может обозначать разные концепции в разных доменах (например, «ядро» в биологии и в операционных системах), что требует контекстного разрешения при нормализации.

\subsection{Анализ решаемого класса задач}
% Определяется класс задач, решаемых решателем: диагностика, классификация, выбор, конфигурирование, планирование, прогнозирование, объяснение или иной класс. Описываются условия задачи, требуемый результат и критерии корректности решения.

Разрабатываемый решатель задач относится к классу информационно-аналитических и объясняющих систем (Explanation and Information Retrieval Systems) с элементами автоматического структурного синтеза баз знаний.

Условия задачи определяются следующим образом:
\begin{itemize}
    \item на вход системы подается либо новый неструктурированный документ, либо запрос пользователя на естественном языке $Q$;
    \item в качестве априорных знаний выступает текущее состояние графа $G$, канонический реестр концепций и векторные пространства эмбеддингов.
\end{itemize}

Требуемый результат зависит от сценария использования:

    \textbf{1 В сценарии пополнения знаний:} корректная интеграция нового документа в существующий граф --- извлечение структурных элементов, нормализация синонимов, построение логических связей с ранее загруженными концепциями и актуализация иерархии.  
    
   \textbf{2 В сценарии поиска ответов:} формирование точного текстового ответа на запрос $Q$ с одновременным построением **объяснительного пути рассуждения** (траектории навигации по узлам графа с указанием конкретных карточек знаний и первоисточников).

Критериями корректности решения задач выступают:

    \textbf{1 Семантическая полнота и непротиворечивость графа:} отсутствие дубликатов сущностей при сохранении канонических имен (точность нормализации синонимов).
    
    \textbf{2 Локальность вычислительной сложности:} сложность выполнения поиска должна составлять $O(\log N)$ за счет двуступенчатого алгоритма, исключающего полный перебор графа.
    
    \textbf{3 Объяснимость и прослеживаемость (Traceability):} каждый сгенерированный языковой моделью ответ должен строго подкрепляться подграфом доказательности со ссылками на конкретные карточки знаний и оригиналы файлов.

Формально задачу поиска ответа можно определить как отображение
\begin{equation}
f: (Q, G) \rightarrow (A, P),
\end{equation}
где $Q$ --- запрос на естественном языке, $G$ --- текущий граф знаний, $A$ --- текстовый ответ, $P = (v_1, \dots, v_k)$ --- путь в графе, обосновывающий ответ. Задача пополнения знаний определяется как
\begin{equation}
g: (D, G, R) \rightarrow G',
\end{equation}
где $D$ --- документ, $R$ --- канонический реестр концепций, $G'$ --- обновлённый граф знаний.


\subsection{Анализ используемой базы знаний}
% Описывается база знаний, разработанная в 4-м семестре или используемая в проекте: её структура, виды знаний, формализм представления, состав понятий, фактов, правил, отношений и ограничений. Определяется, какие элементы базы знаний необходимы решателю задач.

\subsection{Анализ существующих подходов к решению задач}
% Рассматриваются существующие методы, алгоритмы, программные средства и аналогичные интеллектуальные системы. Для каждого подхода указываются преимущества, недостатки и применимость к выбранной предметной области.

Для решения задач организации баз знаний и смыслового поиска в неструктурированных документах в настоящее время применяются три основных подхода: иерархические менеджеры заметок, классические векторные системы RAG и стандартные системы построения графов знаний (Knowledge Graph RAG).


    \textbf{1 Иерархические и сетевые менеджеры заметок (Obsidian, Notion, Roam Research).}
    Данные системы предоставляют пользователю инструменты для создания двоичных и гипертекстовых связей между заметками в ручном или полуавтоматическом режиме.
    \begin{itemize}
        \item \textit{преимущества:} высокая степень контроля над структурой заметок, простота навигации, низкие требования к вычислительным ресурсам;
        \item \textit{недостатки:} высокая трудоемкость ручного связывания, отсутствие автоматической дедупликации концепций и неспособность автоматически извлекать явные семантические единицы из загружаемых внешних файлов большого объема;
        \item \textit{применимость:} подходят для ведения персональных записей, но неприменимы для масштабной автоматизированной обработки неструктурированной информации.
    \end{itemize}

    \textbf{2 Классические векторные системы RAG.}
    Основаны на разбиении текста на фрагменты, получении их векторных представлений и поиске $K$ ближайших фрагментов по метрике косинусного расстояния с последующей передачей контекста в большую языковую модель (LLM).
    \begin{itemize}
        \item \textit{преимущества:} высокая скорость поиска локальных контекстов, простота индексации документов;
        \item \textit{недостатки:} утрата глобального контекста документа, невозможность прослеживания многошаговых логических связей между изолированными фрагментами, отсутствие структурированного представления базы знаний;
        \item \textit{применимость:} эффективны для точечных вопросов по конкретным документам, но не позволяют строить целостную картинку предметной области.
    \end{itemize}

    \textbf{3 Полносвязные графовые системы RAG (Full GraphRAG).}
    Строят масштабные графы сущностей и связей, выполняя глобальный обход или иерархическую суммаризацию всех узлов графа при каждом запросе.
    \begin{itemize}
        \item \textit{преимущества:} глубокий учет семантических связей, высокое качество обобщающих ответов;
        \item \textit{недостатки:} вычислительная сложность уровня $O(V+E)$ на каждый запрос при полном обходе графа, высокая стоимость эксплуатации и задержка ответа при росте базы знаний;
        \item \textit{применимость:} подходят для статических и небольших графов, но требуют оптимизации алгоритмов обхода для работы в реальном времени.
    \end{itemize}


Таким образом, ни один из существующих подходов в чистом виде не удовлетворяет в полной мере требованиям к созданию динамической базы знаний с логарифмической сложностью поиска и прозрачной объяснимостью результатов. Это определяет необходимость разработки комбинированного гибридного подхода на основе локального графового RAG с мультиагентным управлением.


\subsection{Требования к решателю задач}
% Формулируются функциональные требования: поддерживаемые классы задач, входные и выходные данные, использование базы знаний, формирование решения и объяснения.
% Формулируются нефункциональные требования: корректность, полнота, устойчивость, производительность, модульность, расширяемость, сопровождаемость и возможность интеграции с пользовательским интерфейсом.
%
% Данный подраздел завершает аналитическую главу. В нём не следует описывать проектные решения, алгоритмы реализации или программные технологии; формулируются только требования к будущему решателю задач.

На основе проведенного анализа предметной области, существующих подходов и выявленных ограничений сформирован полный комплекс функциональных и нефункциональных требований к проектируемому мультиагентному решателю задач.

\subsubsection{Функциональные требования}


    \textbf{1 Поддерживаемые классы задач и сценарии работы:}
    \begin{itemize}
        \item \textit{автоматический анализ и извлечение знаний:} система должна принимать неструктурированные текстовые документы, определять их тематическую принадлежность, вычленять ключевые тезисы и аннотированное резюме (summary), а также извлекать список базовых предметных концепций;
        \item \textit{каноническая нормализация и дедупликация:} решатель должен контролировать реестр концепций, автоматически выявлять синонимы (в том числе на разных языках или в разных терминологических традициях), сопоставлять новые сущности с имеющимися и предотвращать дублирование узлов графа;
        \item \textit{динамическое построение и обновление связей:} система должна автоматически вычислять и обновлять смысловые и структурные связи (между концепциями, между карточками знаний, а также между концепциями и карточками), рассчитывая вес рёбер на основе частоты совместного упоминания и близости контекстов;
        \item \textit{иерархическая кластеризация знаний:} система должна поддерживать структурирование графа на укрупненные уровни (Сообщество $\rightarrow$ Домен $\rightarrow$ Концепция) с периодическим пересчетом топологии графа и генерацией описаний для верхних уровней;
        \item \textit{интерактивный диалоговый поиск и формирование ответов:} решатель должен принимать вопросы пользователя на естественном языке, искать релевантные зоны знаний и генерировать точные ответы с обязательным построением объяснимого пути рассуждения.
    \end{itemize}

    \textbf{2 Требования к входным данным:}
    \begin{itemize}
        \item поддержка различных форматов неструктурированных документов: файлы формата Portable Document Format (.pdf), текстовые документы Markdown (.md) и файлы Word (.docx);
        \item поддержка текстовых запросов на естественном языке произвольной длины и структуры;
        \item возможность передачи метаданных загружаемого файла (название, дата создания, авторство, путь в файловой системе).
    \end{itemize}

    \textbf{3 Требования к выходным данным:}
    \begin{itemize}
        \item сформированная карточка знания, содержащая аннотацию, основные тезисы, список привязанных концепций и прямую ссылку на исходный документ;
        \item обновленные структуры графа знаний (узлы, рёбра, весовые коэффициенты, канонический реестр синонимов);
        \item текстовый ответ на вопрос пользователя, содержащий логически обоснованные выводы;
        \item графовый путь рассуждения (траектория навигации по узлам вида «вопрос $\rightarrow$ концепция $A$ $\rightarrow$ концепция $B$ $\rightarrow$ первоисточник»), подкрепляющий сгенерированный ответ ссылками на конкретные карточки знаний и оригиналы файлов.
    \end{itemize}

    \textbf{4 Требования к использованию базы знаний и механизму поиска:}
    \begin{itemize}
        \item поиск должен осуществляться по двуступенчатому принципу: первоначальный векторный поиск точек входа по узлам концепций с последующим локальным обходом графа;
        \item локальный обход должен ограничиваться глубиной в 1--2 хопа от найденных точек входа с использованием алгоритмов сжатия фронта поиска (лучевой поиск / \textit{beam search});
        \item система должна производить ранжирование выделенного подграфа по совокупности факторов: векторная близость к запросу, вес рёбер и графовая центральность узлов;
        \item передача контекста в языковую модель для формирования конечного ответа должна ограничиваться строго отфильтрованным локальным подграфом.
    \end{itemize}

    \textbf{5 Требования к формированию объяснения:}
    \begin{itemize}
        \item каждый ответ на запрос пользователя должен сопровождаться явным представлением доказательной базы;
        \item объяснение должно включать цепочку концепций, задействованных при ответе, а также список карточек знаний и оригинальных документов, на основе которых были сделаны выводы;
        \item система должна позволять пользователю перейти от любого элемента в цепочке рассуждения к соответствующей карточке знания и исходному документу.
    \end{itemize}


\subsubsection{Нефункциональные требования}

    \textbf{1 Производительность и вычислительная сложность:}
    \begin{itemize}
        \item вычислительная сложность алгоритма выполнения поискового запроса должна составлять $O(\log N)$ от общего числа узлов базы знаний $N$, исключая полный перебор элементов графа.
        \item время первично поиска точек входа и локального обхода не должно превышать сотен миллисекунд.
        \item суммарное время генерации ответа пользователю с учетом работы языковой модели должно находиться в пределах нескольких секунд.
    \end{itemize}

    \textbf{2 Корректность, полнота и инвариантность знаний:}
    \begin{itemize}
        \item система должна сохранять целостность связей графа при добавлении, обновлении или удалении документов.
        \item ошибки извлечения концепций или некорректный ответ языковой модели не должны приводить к разрушению структуры базы знаний или сбоям в работе канонического реестра.
        \item система должна обеспечивать инвариантность хранения: оригиналы документов хранятся неизменно во внешнем хранилище, а их удаление или перемещение не должно ломать логическую структуру метаданных карточки.
    \end{itemize}

    \textbf{3 Архитектурная модульность и разделение обязанностей:}
    
        Логика системы должна быть строго декомпозирована на автономные программные агенты:
        \begin{itemize}
            \item \textit{агент-аналитик:} ответственность за первичное извлечение структуры из файла;
            \item \textit{агент-нормализатор:} ответственность за реестр канонических концепций и синонимы;
            \item \textit{агент-связыватель:} ответственность за расчет весов и добавление рёбер;
            \item \textit{агент-кластеризатор:} ответственность за регулярное перестроение верхних иерархических уровней;
            \item \textit{агент-ответчик:} ответственность за векторно-графовый поиск и генерацию объяснимых ответов.
        \end{itemize}
        
        Агенты извлечения и нормализации должны работать в фоновом/последовательном режиме при загрузке документа, а агент-ответчик — изолированно в синхронном режиме обработки запроса пользователя.
  

   \textbf{4 Надежность, устойчивость и отказоустойчивость:}
    
        При некорректном формате входного документа или сбое внешней языковой модели система должна корректно изолировать ошибку, уведомлять пользователя и сохранять работоспособность ранее накопленного графа.
        
        Система должна сохранять базовую ценность (структуру графа, реестр концепций и связи между документами) даже в условиях временной недоступности внешней языковой модели.


    \textbf{5 Масштабируемость, расширяемость и сопровождаемость:}
    
        Возможность независимого пересчета топологии графа и сообществ (агентом-кластеризатором) без блокировки или остановки процессов поиска и загрузки новых документов.
        
        Возможность добавления новых типов файлов и парсеров без изменения ядра базы знаний и графового движка.
    

    \textbf{6 Интеграция с пользовательским интерфейсом:}
        
        Решатель должен предоставлять программные интерфейсы (API) для реализации прогрессивного раскрытия информации на клиенте: от верхнего уровня сообществ и доменов к отдельным концепциям, карточкам и первоисточникам.
        
        Поддержка передаче интерфейсу данных о структуре графа в формате, удобном для интерактивной визуализации узлов и рёбер в реальном времени.
        
        Обеспечение отображения прогресса поэтапной обработки загружаемого документа (анализ $\rightarrow$ нормализация $\rightarrow$ связывание).

\subsection{Выводы}

В ходе анализа предметной области установлено, что задача организации личной базы знаний из неструктурированных документов относится к классу информационно-аналитических и объясняющих задач с элементами структурного синтеза. Определены ключевые объекты предметной области (оригинал документа, карточка знания, концепция, канонический реестр, домен, сообщество, семантический граф знаний) и участники процесса.

Анализ существующих подходов показал, что ни один из них в чистом виде не удовлетворяет требованиям динамической базы знаний: иерархические менеджеры заметок требуют ручного связывания; классический RAG утрачивает глобальный контекст и не накапливает знания; полносвязные графовые системы RAG имеют сложность $O(V+E)$ на запрос; онтологический инжиниринг трудоёмок и плохо масштабируется на слабоструктурированные тексты; графовые эмбеддинги не интерпретируются и не учитывают содержание документов.

На основании выявленных ограничений обоснована необходимость гибридного подхода, сочетающего локальный графовый RAG с мультиагентной архитектурой. Сформулированы функциональные и нефункциональные требования к решателю задач, включая требования к производительности, модульности, надёжности и интеграции с пользовательским интерфейсом.

Сформированные требования являются исходными данными для проектирования архитектуры системы, рассматриваемой в следующей главе.
